\documentclass[11pt]{article}
\usepackage[margin=1in]{geometry}
\usepackage{amsmath,amssymb,amsthm,mathrsfs,mathtools,microtype}
\usepackage[hidelinks]{hyperref}

\newtheorem{theorem}{Theorem}[section]
\newtheorem{proposition}[theorem]{Proposition}
\newtheorem{lemma}[theorem]{Lemma}
\newtheorem{corollary}[theorem]{Corollary}
\newtheorem{remark}[theorem]{Remark}
\newcommand{\Mbar}{\overline{M}}
\newcommand{\Norm}{\operatorname{Norm}}
\newcommand{\Spec}{\operatorname{Spec}}

\title{The Missing H1/H2 Comparison Package:\\
An Obstruction and the Correct Common Resolution}
\author{Repository comparison note}
\date{23 July 2026}

\begin{document}
\maketitle

\begin{abstract}
We test the proposed comparison
\[
 \mathscr S_N^{\mathrm{seed}}
 \longrightarrow \mathscr H_N^{\mathrm{adm}}
 \longrightarrow \mathcal C_N
\]
one arrow at a time.  The first arrow does not exist for the compactifications
currently assigned these names.  Already in degree five, the stable
zero-root configuration forgets a ratio of branch-value scales that the
fully branch-marked admissible target retains.  Locally the missing datum is
the rational function $x^3/y^2$, and resolving it requires the normalized
blowup of $(x^3,y^2)$.  Thus the admissible-cover closure is a modification,
not an identified copy, of the root-stable quotient.

We replace the false arrow by the normalization of the graph.  The second
arrow, from the normal admissible incidence to the repository's normalized
finite incidence, is then constructed by the universal property of
normalization.  We state exact criteria for finiteness, birationality,
normalization identification, compatibility with the selected affine root,
and agreement of completed local rings.  This settles the comparison
question negatively as stated and gives the minimal corrected package needed
by downstream arguments.
\end{abstract}

\section{The objects and the claimed arrows}

Work over an algebraically closed field $k$ of characteristic zero.  Put
$n=N-2$.  Let
\[
 \mathscr S_N^{\mathrm{seed}}
   =[\Mbar_{0,N}/S_{n-1}]
\]
be the collision-separating compactification with the simple root $1$
selected, and let $\mathscr S_N^\circ$ be its smooth, distinct-root open.
Let $\mathscr H_N^{\mathrm{adm}}$ be the normalization of the closure of
this open in a separated admissible-cover stack which retains the stabilized
target branch divisor.  All finite branch points may be unordered; the
argument below only uses whether their stabilized target is reducible.

Finally, let $B_N$ be the chosen repository base and let $\mathcal C_N$ be
the normalization of the scheme-theoretic closure of the relevant generic
finite incidence over $B_N$.  Thus $\mathcal C_N\to B_N$ is finite and
normal and comes with a fixed generic incidence algebra.  This definition
includes either the residual degree-$n$ root cover or, after replacing the
generic algebra, the degree-$N$ inverse incidence.  The two degrees must not
be mixed.

On $\mathscr S_N^\circ$ the polynomial and hence its admissible cover are
defined.  This gives a rational map
\[
 \phi^\circ:\mathscr S_N^{\mathrm{seed}}
        \dashrightarrow\mathscr H_N^{\mathrm{adm}}.       \tag{1.1}
\]
The issue is whether (1.1) is regular at the boundary.

\section{A degree-five obstruction to the first arrow}

Take $N=5$ and write the simple zero-fiber roots as $1,a,b$.  On the clean
open put
\[
 F_{a,b}(W)=W^2(W-1)(W-a)(W-b),\qquad
 H_{a,b}(W)=-\frac{F_{a,b}(W)}{F'_{a,b}(1)}.          \tag{2.1}
\]
Then $H(0)=H'(0)=H(1)=0$ and $H'(1)=-1$.

Near the codimension-two boundary point at which $a$ collides with $0$ and
$b$ collides with $1$, write
\[
 a=x,\qquad b=1+y.                                   \tag{2.2}
\]
The completed local ring of the labelled root-stable space at this point is
$k[[x,y]]$.  The special stable source is independent of the relative rates
of $x$ and $y$: it has a $0$-tail carrying $(0,a)$, a $1$-tail carrying
$(1,b)$, and a central component carrying $\infty$ and the two attaching
nodes.

There are two moving critical values near the marked target value zero.
One comes from the triple cluster $(0,0,a)$ and the other from the pair
$(1,b)$.  Denote them by $\lambda_0$ and $\lambda_1$, respectively.

\begin{lemma}[branch-scale asymptotics]\label{lem:asymptotics}
In the fraction field of $k[[x,y]]$, along every trait on which
$3v(x)>v(y)>0$, one has
\[
 \lambda_0=-\frac4{27}\frac{x^3}{y}\,(1+o(1)),
 \qquad
 \lambda_1=-\frac14y\,(1+o(1)).                     \tag{2.3}
\]
Consequently
\[
 \frac{\lambda_0}{\lambda_1}
   =\frac{16}{27}\frac{x^3}{y^2}\,(1+o(1)).         \tag{2.4}
\]
\end{lemma}

\begin{proof}
Since $F'_{a,b}(1)=(1-a)(1-b)=-(1-x)y$, equation (2.1) gives
\[
 H_{a,b}(W)=\frac{F_{a,b}(W)}{(1-x)y}.                \tag{2.5}
\]
Put $W=xX$.  Uniformly along the stated traits,
\[
 H_{a,b}(xX)=\frac{x^3}{y}X^2(X-1)(1+o(1)).
\]
Besides the fixed critical point $X=0$, the moving critical point tends to
$X=2/3$, giving the first formula in (2.3).  Put instead $W=1+yX$.  Then
\[
 H_{a,b}(1+yX)=yX(X-1)(1+o(1)),
\]
whose moving critical point tends to $X=1/2$.  This gives the second formula
and hence (2.4).
\end{proof}

\begin{theorem}[nonextension]\label{thm:no-arrow}
For $N\geq5$, the rational map (1.1) does not extend to a morphism from the
root-stable quotient to any separated admissible-cover compactification that
retains the stabilized target branch divisor.
\end{theorem}

\begin{proof}
It is enough to work on the labelled degree-five chart; fixed spectator roots
give the assertion in every higher degree.  Consider two traits through the
same point $x=y=0$:
\[
 \begin{array}{c|cc|c}
 &x&y&\lambda_0/\lambda_1\\ \hline
 T_A&t&t& (16/27)t(1+o(1))\\
 T_B&t^2&t^3& (16/27)(1+o(1)).
 \end{array}                                          \tag{2.6}
\]
Both have the same stable marked source described after (2.2).  On the
target, retain the four points $0,\lambda_0,\lambda_1,\infty$ and forget all
other branch points.  Along $T_A$ their stable four-pointed limit lies on the
boundary of $\Mbar_{0,4}$ because $\lambda_0/\lambda_1\to0$.  Along $T_B$
the limit is smooth because the ratio tends to $16/27\notin\{0,1,\infty\}$.
The same distinction survives if $\lambda_0$ and $\lambda_1$ are unordered.

Thus the two admissible covers have nonisomorphic stabilized targets.  If a
morphism (1.1) existed, both traits would send their common closed source
point to one point of the admissible-cover stack.  Separatedness and stable
reduction force that point to be each of the two nonisomorphic limits, a
contradiction.
\end{proof}

\begin{corollary}[the missing local modification]\label{cor:blowup}
At the boundary point above, the rational target-stabilization map contains
the rational map
\[
 \Spec k[[x,y]]\dashrightarrow\mathbb P^1,
 \qquad [x^3:y^2].                                   \tag{2.7}
\]
Its graph is the blowup of $(x^3,y^2)$, and its normal graph is the
normalized blowup of this monomial ideal.  In particular the graph
projection is not finite and the completed local ring is not the original
$k[[x,y]]$.
\end{corollary}

\begin{proof}
Equation (2.4) identifies the target cross-ratio up to a unit with
$x^3/y^2$.  The closure of the graph of a rational map $[f:g]$ is the blowup
of $(f,g)$.  Its fiber over $(x,y)=(0,0)$ contains the exceptional projective
curve; normalization does not make that fiber finite.
\end{proof}

This is the precise point missed by a proof based only on DVR extension.
Every individual trait has an admissible stable limit, but those limits do
not descend to a morphism from the unmodified two-dimensional base.

\section{The corrected first comparison}

Let
\[
 \Gamma_N=\Norm\overline{\Gamma_{\phi^\circ}}
 \subset
 \mathscr S_N^{\mathrm{seed}}\times
 \mathscr H_N^{\mathrm{adm}}                       \tag{3.1}
\]
be the normalization of the graph closure, interpreted on labelled atlases
and descended equivariantly to the quotient stacks.  It comes with
\[
 \mathscr S_N^{\mathrm{seed}}
 \xleftarrow{\ p\ }\Gamma_N
 \xrightarrow{\ q\ }\mathscr H_N^{\mathrm{adm}}.   \tag{3.2}
\]

\begin{proposition}[common-resolution package]\label{prop:graph}
The maps in (3.2) are proper and birational onto the components containing
$\mathscr S_N^\circ$.  The following are equivalent:
\begin{enumerate}
\item $\phi^\circ$ extends to
$\mathscr S_N^{\mathrm{seed}}\to\mathscr H_N^{\mathrm{adm}}$;
\item $p$ is an isomorphism;
\item $p$ is finite.
\end{enumerate}
For $N\ge5$ these conditions fail for the fully branch-marked admissible
compactification.
\end{proposition}

\begin{proof}
Properness follows from graph closure in the product of proper stacks, and
birationality follows because both projections restrict to the identity on
the common dense open.  An extension gives a section of $p$ whose graph is
all of $\Gamma_N$; separatedness gives uniqueness, hence $p$ is an
isomorphism.  Conversely an inverse to $p$ gives the extension.  If $p$ is
finite, it is finite birational onto the normal stack
$\mathscr S_N^{\mathrm{seed}}$, hence is an isomorphism.  Theorem
\ref{thm:no-arrow} and Corollary~\ref{cor:blowup} give the final assertion.
\end{proof}

Thus the corrected first arrow is not
$\mathscr S_N^{\mathrm{seed}}\to\mathscr H_N^{\mathrm{adm}}$ but the
correspondence (3.2).  Equivalently, one may replace the source by the
logarithmic/toroidal modification $\Gamma_N$ that records relative
branch-value scales.  In the local chart (2.2), the first required
subdivision is the normalized blowup of $(x^3,y^2)$.

\section{Construction of the incidence contraction}

The second comparison arrow is formal once the generic algebra and the base
are fixed correctly.

\begin{proposition}[normalization contraction]\label{prop:contraction}
Let $B$ be an integral locally Noetherian base, let $K/k(B)$ be the fixed
finite generic incidence algebra (a field component, or a specified product
treated componentwise), and let
\[
 C=\Norm_B(K)
\]
be finite over $B$.  Let $H$ be a normal integral algebraic stack, dominant
and proper over $B$, whose generic function field is the chosen component of
$K$.  Then the generic comparison extends uniquely to a $B$-morphism
\[
 c:H\longrightarrow C.                              \tag{4.1}
\]
It is proper and birational.  It is finite if and only if it is
quasi-finite, and in that case it is an isomorphism.
\end{proposition}

\begin{proof}
The assertion is smooth-local on $H$ and affine-local on $B$.  If
$B=\Spec A$ and $C=\Spec\overline A$, every element of $\overline A$ is
integral over $A$.  Its image in the common function field is therefore
integral over every normal local ring of $H$.  It belongs to that local ring,
so the generic map extends and the local extensions glue uniquely.  This is
the universal property of normalization.  Properness follows because $H$ is
proper over $B$ and $C$ is separated over $B$.  The generic comparison makes
$c$ birational.  A proper quasi-finite morphism is finite; a finite
birational morphism to normal $C$ is an isomorphism.
\end{proof}

Apply Proposition~\ref{prop:contraction} with
$H=\mathscr H_N^{\mathrm{adm}}$ and $C=\mathcal C_N$, separately for:
\begin{enumerate}
\item the residual degree-$n$ zero-fiber algebra after removing the fixed
length-two summand; and
\item the degree-$N$ inverse-incidence algebra
$k(s,t)[W]/(H(W)-sW+t)$.
\end{enumerate}
This constructs the arrow
\[
 \mathscr H_N^{\mathrm{adm}}\longrightarrow\mathcal C_N              \tag{4.2}
\]
provided the admissible object in (4.2) is the corresponding marked-point
incidence and its generic algebra has first been verified component by
component.  Normalized Stein factorization is an equivalent construction:
the normal $H$ maps to the normalization of
$\Spec_B(f_*\mathcal O_H)$ in the fixed generic algebra.

\section{Compatibility with the selected affine root}

On the clean open, the chosen point $p_1$ maps under (4.2) to the
tautological root of the finite incidence.  Both $\mathcal C_N\to B_N$ and
the coarse marked root cover are separated.  Therefore equality on the dense
open implies equality after pullback to $\Gamma_N$ wherever the two sections
extend.  The finite-cover DVR lemma then gives, for each fixed generic
selected root, a unique extension over every trait.

This proves traitwise compatibility but not descent from $\Gamma_N$ to
$\mathscr S_N^{\mathrm{seed}}$.  Descent would require the selected section
to be constant on every fiber of $p$ in (3.2).  The abstract selected-root
cover does satisfy that statement, but the full admissible target does not:
Corollary~\ref{cor:blowup} exhibits its varying exceptional parameter.

\section{Completed local rings and the exact final status}

Let $\gamma\in\Gamma_N$ map to $h\in\mathscr H_N^{\mathrm{adm}}$ and
$z\in\mathcal C_N$.  The morphism (4.2) induces
\[
 \widehat{\mathcal O}_{\mathcal C_N,z}
 \longrightarrow
 \widehat{\mathcal O}_{\mathscr H_N^{\mathrm{adm}},h}.              \tag{6.1}
\]
An isomorphism in (6.1) is a consequence of an actual global morphism plus a
formal calculation; it cannot be used to manufacture that morphism.
For a proper birational map of excellent normal spaces, isomorphism of these
completed maps at every point implies that the map is formally
quasi-finite, hence quasi-finite, finite, and finally an isomorphism.

At a pure root collision with all target-scale data fixed, the repository's
invariant calculation
\[
 k[[e_1,\ldots,e_\mu,T]]/
 (T^\mu-e_1T^{\mu-1}+\cdots+(-1)^\mu e_\mu)          \tag{6.2}
\]
is the correct completed finite root cover.  It proves agreement of the
root-selection factor.  It does not include the target cross-ratio
$x^3/y^2$ in (2.4).  At the simultaneous collision of Section~2 the source
completion is $k[[x,y]]$, whereas the graph/admissible side contains the
normalized blowup chart of $(x^3,y^2)$.  Hence the full completed local rings
do not agree.

The requested checklist therefore has the following unconditional answer:
\begin{center}
\begin{tabular}{ll}
boundary extension & yes after replacing the first arrow by $\Gamma_N$;\\
normalization compatibility & yes, by Proposition~\ref{prop:contraction};\\
selected affine root & yes on $\Gamma_N$ and on every trait;\\
finiteness of $\Gamma_N\to\mathscr S_N^{\rm seed}$ & no for $N\ge5$;\\
birationality onto the claimed components & yes;\\
normality/normalization identification & yes for $\mathcal C_N$ by definition;\\
full completed-ring agreement & no at the two-scale boundary;\\
root-factor completed-ring agreement & yes, by (6.2).
\end{tabular}
\end{center}

Accordingly, H1/H2 cannot be repaired by proving the originally displayed
arrows.  The correct global package is
\[
 \boxed{
 \mathscr S_N^{\mathrm{seed}}
 \xleftarrow{\text{proper birational}}
 \Gamma_N
 \xrightarrow{\text{proper birational}}
 \mathscr H_N^{\mathrm{adm}}
 \xrightarrow{\text{normalization contraction}}
 \mathcal C_N .}
\]
Downstream statements using only the finite selected-root cover and its DVR
extension survive.  Statements identifying the entire root-stable and
branch-stable compactifications must either use $\Gamma_N$ or forget enough
target branch-scale data for the obstruction to disappear.

\end{document}
